% =============================================================================
% Midterm cheatsheet, Topics 1, 2A, 2B (Finanças II, Prof. Ricca).
% Concepts, solving recipes, "watch out" lists, explain-templates and a True/False bank.
% Complements Notas/formula-sheet-topics-1-2b.pdf (formulas) -- this sheet is about HOW to solve.
% Compile: pdflatex (twice); check: pdffonts (Type 1 only, uni = yes).
% =============================================================================
\documentclass[10pt,a4paper]{article}
\usepackage[utf8]{inputenc}
\usepackage[T1]{fontenc}
\usepackage{lmodern}   % Type 1 vector fonts (replaces PK bitmaps)
\usepackage{cmap}      % ToUnicode CMap -> copyable text and accents
\usepackage[english]{babel}
\usepackage{amsmath,amssymb}
\usepackage[a4paper,margin=1.7cm]{geometry}
\usepackage{enumitem}
\usepackage{xcolor}
\usepackage{fancyhdr}
\usepackage{booktabs}
\usepackage{tabularx}
\usepackage{multicol}
\usepackage{microtype}

\setlist{nosep,leftmargin=1.3em}
\pagestyle{fancy}\fancyhf{}
\lhead{\footnotesize Finan\c{c}as II --- Corporativas (MPE Insper)}
\rhead{\footnotesize Midterm cheatsheet --- Topics 1, 2A, 2B}
\cfoot{\footnotesize\thepage}
\renewcommand{\headrulewidth}{0.3pt}
\setlength{\headheight}{14pt}
\setlength{\parindent}{0pt}
\setlength{\parskip}{3pt}
\setlength{\columnsep}{14pt}
\newcolumntype{Y}{>{\raggedright\arraybackslash}X}

\definecolor{hd}{RGB}{25,60,120}
\definecolor{wr}{RGB}{160,30,30}
\definecolor{ok}{RGB}{20,110,40}
\newcommand{\blk}[1]{\subsection*{\color{hd}#1}}
\newcommand{\recipe}[1]{\par\smallskip\noindent\textbf{\color{hd}Recipe --- #1}\par}
\newcommand{\watch}{\par\noindent\textbf{\color{wr}Watch out:}~}
\newcommand{\T}{\textbf{\color{ok}[T]}~}
\newcommand{\F}{\textbf{\color{wr}[F]}~}
\newcommand{\E}{\mathbb{E}}

\begin{document}

\begin{center}
  {\LARGE\bfseries Midterm cheatsheet --- Topics 1, 2A and 2B}\\[3pt]
  {\large Concepts, how to solve, what to watch, and a True/False bank}
\end{center}
\vspace{-4pt}
\noindent\textbf{How this sheet fits.} Formulas are in \texttt{Notas/formula-sheet-topics-1-2b.pdf}; this
sheet is about \emph{using} them: what each concept means, step-by-step recipes for the question types
seen in last year's midterm (\texttt{Prova/q2.pdf}), PS1--PS2 and the mock
(\texttt{Simulados/mock-midterm-topics-1-2b.pdf}), the traps, ready-made ``explain'' answers, and 75
True/False statements with the reason.

% =============================================================================
\section*{0. Exam-day checklist (every question)}
\begin{multicols}{2}
\begin{enumerate}[label=\arabic*.]
  \item \textbf{Units and scale}: thousands vs.\ millions; per share vs.\ total; \% vs.\ decimal
        (write $0.20$, not $20$, inside formulas).
  \item \textbf{Timing}: is the first cash flow at $t=0$ or $t=1$? A perpetuity $C/r$ starts
        \emph{next} period; a growing perpetuity uses $C_1$, not $C_0$.
  \item \textbf{Annuity vs.\ perpetuity}: ``for 5 years'' $\Rightarrow$ annuity factor; ``each year''
        with no end $\Rightarrow$ perpetuity.
  \item \textbf{Which world?} No taxes / corporate taxes / personal taxes / distress costs / risky vs.\
        risk-free debt. Write it down before choosing a formula.
  \item \textbf{Which debt policy?} Fixed amount (permanent), constant $D/E$, or a known schedule.
        It decides $r_{ITS}$ and the formula for $r_E$ (table in §3).
  \item \textbf{Book or market?} Ratios, weights in WACC and betas use \emph{market} values.
  \item \textbf{State the formula, then plug in}. Partial credit lives in the formula line.
  \item \textbf{Keep 4 decimals} in intermediate steps; round only the final answer.
  \item \textbf{Sanity checks}: $r_D\le r_U\le r_E$; $r_D(1-\tau)<r_{WACC}<r_E$; $r_{WACC}\le r_U$
        with taxes; $V_L\ge V_U$ without distress; betas: $\beta_D\le\beta_A\le\beta_E$.
  \item \textbf{Answer every sub-question}: ``Are investors willing?'', ``smaller, equal or larger?'',
        ``Explain'': each is worth marks. Use one sentence of mechanism, not a definition.
  \item \textbf{Cross-check with a second route} when cheap: WACC vs.\ APV; CAPM vs.\ MM II; FCF from
        EBIT vs.\ from net income.
  \item \textbf{Interpret the sign}: negative NPV $\Rightarrow$ not financed, not undertaken.
\end{enumerate}
\end{multicols}

% =============================================================================
\section*{1. Topic 1 --- Accounting review and the capital structure}

\blk{Core concepts (one line each)}
\begin{itemize}
  \item \textbf{Balance sheet} $A=L+E$ holds in book \emph{and} market value. Liabilities are
        promises (senior, fixed); equity is the residual (junior). Firms sell goods \emph{and}
        promises; the promise side is the capital structure.
  \item \textbf{Book vs.\ market}: book is historical cost, backward-looking, omits most intangibles
        (about 90\% of S\&P 500 EV); market is forward-looking. Book equity can be negative while
        market equity is large (Starbucks 2020: $-7.9$ bn book vs.\ $122.6$ bn market cap).
        SVB: held-to-maturity bonds at book hid market losses.
  \item \textbf{Enterprise value} $EV=E^{mkt}+D-\text{Cash}$: value of the operating business,
        owned jointly by equity and net debt. Pair it with firm-level flows (sales, EBITDA, FCF).
  \item \textbf{Accrual vs.\ cash}: revenue when earned, expense when incurred (matching).
        Depreciation spreads capex over time and is \emph{not} a cash outflow.
  \item \textbf{NWC} $=CA-CL$ (operating items): money tied up in running the business; growth
        absorbs cash. Payables are free short-term finance.
  \item \textbf{Operating leverage}: high fixed costs $\Rightarrow$ $DOL=1+F/OI$ large $\Rightarrow$
        operating income swings more with sales.
  \item \textbf{ROE vs.\ ROIC}: $ROE=\frac{NI}{S}\cdot\frac SA\cdot\frac AE$ mixes operations and
        financing; $ROIC=EBIT(1-\tau)/(D+E)$ measures operations only.
        $ROE=ROIC+\frac DE[ROIC-r_D(1-\tau)]$: leverage helps ROE only if $ROIC>r_D(1-\tau)$, and
        always makes it riskier.
  \item \textbf{FCF vs.\ FCFE}: FCF (unlevered) goes to all capital providers $\to$ discount at WACC
        (or $r_U$ in APV) $\to$ EV. FCFE goes to equity $\to$ discount at $r_E$ $\to$ equity value.
\end{itemize}

\recipe{Free cash flow from an income statement or a consultant's report (PS1 Ex.\,2, Berk 9.13, 9.25)}
\begin{enumerate}
  \item Strip the statement down to \emph{incremental} items: drop sunk costs (the consultant's fee,
        past R\&D), drop allocated overhead that occurs anyway; add opportunity costs and
        cannibalisation.
  \item $EBIT=\text{Revenue}-\text{COGS}-\text{incremental Opex}-\text{Dep}$; never subtract interest.
  \item Unlevered net income $=EBIT(1-\tau)$.
  \item $FCF=EBIT(1-\tau)+\text{Dep}-\text{Capex}-\Delta NWC$. Equivalent:
        $(\text{Rev}-\text{Costs})(1-\tau)+\tau\,\text{Dep}-\text{Capex}-\Delta NWC$.
  \item NWC: invest it when needed (often $t=0$), recover it at the end (a positive flow).
  \item Salvage: after-tax $=SV-\tau(SV-BV)$.
\end{enumerate}
\watch depreciation enters FCF only through its tax shield $\tau\,\text{Dep}$; if the report already
deducted interest, add $\text{Int}(1-\tau)$ back; an increase in NWC is a cash \emph{outflow}; a
\emph{decrease} is an inflow; if the tax rate changes by year, use each year's rate.

\recipe{Cash flow from operations (indirect method)}
$CFO=NI+\text{Dep}+\text{other non-cash}-\Delta AR-\Delta\text{Inv}+\Delta AP$ (+ increase in deferred
tax liabilities, PS1 Ex.\,1b). $\Delta\text{Cash}=CFO+CFI+CFF$.
\watch signs: assets up $\Rightarrow$ cash down; liabilities up $\Rightarrow$ cash up. Capex is in CFI,
dividends and borrowing in CFF.

\recipe{Multiples and leverage ratios (Berk 2.8, 2.21, 2.24)}
Compute $E^{mkt}=P\times N$ first, then $EV$, then the ratio. State whether $D/E$ is book or market.
\watch price-to-sales and P/E are distorted by leverage; EV/sales and EV/EBITDA are not. Interest
coverage $=EBIT/\text{Int}$ measures the ability to pay, which leverage ratios do not.

% =============================================================================
\section*{2. Topic 2A --- Valuation, NPV and the frictionless benchmark}

\blk{Core concepts}
\begin{itemize}
  \item \textbf{Price $=$ expected cash flows discounted at the expected return} that investors require
        for that risk: $r=r_f+\text{risk premium}$. Expected $\ne$ realised return.
  \item \textbf{The discount rate belongs to the project}, not to the firm or to whoever finances it.
        Get it from \emph{comparables}: unlever each comparable's $\beta_E$, average the $\beta_A$.
  \item \textbf{Unlevering}: returns and betas are value-weighted averages,
        $\beta_A=\frac EV\beta_E+\frac DV\beta_D$; with safe debt $\beta_A=\frac EV\beta_E$.
        Slide example: a house bought with 20\% or 50\% debt has $\beta_E=1.25$ or $2$ with the same
        $\beta_A=1$.
  \item \textbf{NPV rule}: accept iff $NPV=PV-\text{cost}>0$. In a frictionless market the NPV goes to
        the existing owners \emph{whatever the financing} (internal funds, new equity, safe debt).
        Positive NPV is ``arbitrage across the firm's boundary'': it cannot be scaled, so it survives.
  \item \textbf{Fisher separation}: investment is chosen at the market rate; owners' patience does not
        matter (they borrow or lend on their own).
  \item \textbf{IRR pitfalls}: multiple or no IRR (sign changes), borrowing-type flows (accept if
        $IRR<r$), scale and timing in mutually exclusive projects. Use NPV.
  \item \textbf{Debt and equity as options}: $D=\min(A,F)=$ safe bond $-$ put; $E=\max(A-F,0)=$ call
        with strike $F$. Put--call parity $\Rightarrow$ $D+E=A$ for any $F$: that is MM I.
  \item \textbf{Complete markets}: state prices $p_s$ price everything; $\sum_s p_s=1/(1+r_f)$; no
        arbitrage $\Leftrightarrow$ all $p_s>0$. Slump claims cost more than their probability
        suggests (risk aversion). Risk-neutral probabilities $\pi_s=p_s(1+r_f)$ are not real ones.
  \item \textbf{MM I}: $V_L=V_U$. Proof by arbitrage: homemade leverage (borrow on personal account)
        or undo leverage (buy the firm's debt). Needs: same borrowing terms, no taxes, no distress
        costs, no information or agency effects, fixed investment.
  \item \textbf{MM II}: $r_E=r_A+\frac DE(r_A-r_D)$; WACC $=r_A$ is unaffected by leverage. Higher
        $r_E$ is pay for higher risk, not a gain. ``Conservation of value'': financial transactions
        repackage risk.
  \item \textbf{$r_D$ is not the yield}: with default risk, the promised yield $y$ exceeds the expected
        return $r_D$.
\end{itemize}

\recipe{Project NPV}
(1) FCF by year (§1 recipe) $\to$ (2) discount rate from comparables (below) $\to$ (3) $NPV=\sum
FCF_t/(1+r)^t$, with an annuity or perpetuity where possible $\to$ (4) decide and explain.
\watch the investment at $t=0$ is negative; NWC recovered in the last year; terminal value
$TV_T=FCF_{T+1}/(r-g)$ is discounted $T$ periods, not $T+1$; real with real, nominal with nominal.

\recipe{Cost of capital from comparables (Prova Q2(a), Lista R2 ex.\,6)}
\begin{enumerate}
  \item For each comparable: market weights $E/V$, $D/V$ (from $D/E$: $D/V=\frac{D/E}{1+D/E}$).
  \item $\beta_A=\frac EV\beta_E+\frac DV\beta_D$ ($\beta_D=0$ if debt is risk-free). With taxes and
        constant $D/E$, $\beta_A=\beta_U$ (the shields have the risk of $U$).
  \item Average the $\beta_A$; $r_U=r_f+\beta_U(\E[r_m]-r_f)$.
  \item Relever at the \emph{project's} target $D/E$ if the WACC method is used.
\end{enumerate}
\watch never use a comparable's $\beta_E$ directly; never average equity betas of firms with different
leverage; the \emph{excess} market return $\E[r_m]-r_f$ may be given instead of $\E[r_m]$; read which.

\recipe{State prices, options and MM in a complete market (PS2 Ex.\,1)}
\begin{enumerate}
  \item Two states, two traded assets: write $\text{price}=\sum_s p_s x_s$ for each; solve the
        $2\times2$ system.
  \item Price any claim by $\sum_s p_sx_s$ (calls, puts, the project, debt with face $F$).
  \item Debt of face $F$: payoff $\min(X_s,F)$; equity $\max(X_s-F,0)$; check $D+E=A$.
  \item Expected returns use the \emph{true} probabilities: $r=\E[\text{payoff}]/\text{price}-1$.
  \item To raise an amount $D$: find $F$ such that $\sum_s p_s\min(X_s,F)=D$; the case where $F$
        exceeds the slump payoff makes debt risky.
\end{enumerate}
\watch state prices are not probabilities; $r_A$ is fixed by the project and does not move with $D$;
a debt-financed dividend does not raise firm value: equity falls by the dividend (PS2 1f).

\recipe{Leverage change in a perfect market: $r_E$, EPS, P/E (PS2 Ex.\,4, Berk 16.5)}
Unlever ($r_A$, $\beta_A$) $\to$ relever at the new $D/E$: $r_E=r_A+\frac DE(r_A-r_D)$,
$\beta_E=\beta_A+\frac DE(\beta_A-\beta_D)$. EPS after $=(\text{EBIT}-r_DD)/N_{new}$. Price per share
is \emph{unchanged} (MM), so P/E $=P/EPS_{new}$.
\watch EPS can rise with leverage (if the earnings yield exceeds $r_D$), but the price does not: the
higher EPS is riskier. A falling P/E is then correct: it reflects the higher risk, not cheapness. ``The
capital structure that maximises the expected return on equity is best'' is wrong.

% =============================================================================
\section*{3. Topic 2B --- Taxes, distress and the static trade-off}

\blk{Core concepts}
\begin{itemize}
  \item \textbf{Interest tax shield}: interest is deductible, dividends are not $\Rightarrow$ the
        government takes less from a levered firm. $ITS_t=\tau r_D D_{t-1}$.
        $V_L=V_U+PV(ITS)$. The gain goes to \emph{shareholders} (debt is fairly priced).
  \item \textbf{Unlevered FCF ignores the shield}; it enters through the after-tax WACC or through
        APV, never both.
  \item \textbf{Risk of the shield depends on the debt policy} (table below). This is the single most
        important ``which formula?'' decision in 2B.
  \item \textbf{100\%-debt puzzle}: corporate taxes alone imply all-debt firms; observed leverage is
        moderate. Answers: personal taxes (Miller), distress costs, agency costs.
  \item \textbf{Default vs.\ bankruptcy costs}: default is a transfer of control and does not break MM;
        the \emph{costs} of default do. Direct (legal, administrative, fire sales) and indirect (lost
        customers, suppliers cutting credit, workers leaving, debt overhang, risk-shifting), which
        start \emph{before} any filing (Hortaçsu et al.\ 2013: used-car prices and CDS).
  \item \textbf{Who pays distress costs}: shareholders, ex ante. Creditors anticipate them and lend
        less for a given face value, so equity raises less (PS2 Ex.\,3).
  \item \textbf{Static trade-off}: choose debt where the marginal tax benefit equals the marginal
        expected distress cost: $\tau[1-H(K^*)]=C\,h(K^*)$. $K^*$ rises with $\tau$ and falls with $C$.
        Predicts more debt for stable, tangible, profitable firms; less for R\&D, brand and
        service-dependent firms.
  \item \textbf{Bank capital (Admati et al.)}: more equity lowers expected ROE \emph{and} its risk;
        required $r_E$ falls and the WACC is unchanged (MM II). The only private cost is the lost tax
        shield, a transfer from taxpayers, not a social cost. Failures impose an externality, so the
        social optimum has more equity than banks choose.
  \item \textbf{Payout irrelevance (MM)}: with investment fixed, a higher dividend must be financed by
        issuing claims of equal value; dividend vs.\ repurchase does not matter in a perfect market.
        Taxes break it: dividends taxed more than capital gains favour repurchases; corporate tax on
        interest makes retained excess cash costly (Berk 17.12).
\end{itemize}

\blk{Which formula? The debt-policy table (the 2B decision)}
\begin{center}\small
\begin{tabularx}{\linewidth}{@{}>{\raggedright\arraybackslash}p{3.4cm}>{\raggedright\arraybackslash}p{3.6cm}Y>{\raggedright\arraybackslash}p{2.5cm}@{}}
  \toprule
  Debt policy & $r_{ITS}$, $PV(ITS)$ & $r_E$ and $\beta_E$ & WACC shortcut \\ \midrule
  Permanent fixed amount $D$ (``maintain a balance of'') & $r_D$;\ $PV=\tau D$ &
  $r_U+\frac DE(1-\tau)(r_U-r_D)$;\ $\beta_E=\beta_U[1+(1-\tau)\frac DE]$ (safe debt) &
  $r_U(1-\tau D/V)$ ($g=0$) \\[2pt]
  Constant $D/E$ or $D/V$ (``target leverage'') & $r_U$;\ discount each $ITS_t$ at $r_U$ &
  $r_U+\frac DE(r_U-r_D)$;\ $\beta_E=\beta_U+\frac DE(\beta_U-\beta_D)$ &
  $r_U-\tau r_D\frac DV$ \\[2pt]
  Known schedule (loan repaid over $T$ years) & $r_D$ (shields as safe as the debt);\ sum
  $ITS_t/(1+r_D)^t$ & use APV, not WACC & --- \\ \bottomrule
\end{tabularx}
\end{center}
\watch the word ``maintain a constant debt-to-equity ratio'' $\Rightarrow$ second row (Prova Q2, PS2
Ex.\,2); ``issues \$100 million of perpetual debt'' $\Rightarrow$ first row (PS2 Ex.\,2e).

\recipe{Debt-financed repurchase (Berk 16.18, 16.22, slide example 600/500/40\%)}
\begin{enumerate}
  \item Before: $V_U=P_0N_0$.
  \item Announcement: $V_L=V_U+PV(ITS)$ ($=\tau D$ if permanent); new price
        $P=V_L/N_0$ (the gain is capitalised \emph{now}, before any share is bought).
  \item Repurchase at $P$: shares bought $=D/P$; $N_1=N_0-D/P$.
  \item After: $E=V_L-D$, price $=E/N_1=P$ (unchanged by the repurchase itself).
\end{enumerate}
\watch with distress costs, $V_L=V_U+PV(ITS)-PV(\text{distress})$; if the new price is given, back
out $PV(\text{distress})=V_U+PV(ITS)-P_{new}N_0$ (Berk 16.22).

\recipe{The Prova Q2 template: comparables, constant $D/E$, WACC vs.\ APV}
\begin{enumerate}
  \item Comparable: $\beta_A=\beta_U=\frac EV\beta_E$ (safe debt, constant $D/E$) $\to$ $r_U$ by CAPM.
  \item All-equity: $V_U=\sum FCF_t/(1+r_U)^t$; $NPV_U=V_U-I$. If $NPV<0$, outsiders would need more
        than 100\% of the firm $\Rightarrow$ no financing.
  \item Levered at target $D/E$: $r_E=r_U+\frac DE(r_U-r_D)$; $r_{WACC}=\frac EVr_E+\frac DVr_D(1-\tau)$
        ($=r_U-\tau r_D\frac DV$); $V_L=\sum FCF_t/(1+r_{WACC})^t$.
  \item Balance sheet at 0: $D_0=\frac DVV_L$, $E_0=\frac EVV_L$. Outside equity provides $I-D_0$ and
        gets the fraction $(I-D_0)/E_0$; owners keep the rest, worth $E_0-(I-D_0)=V_L-I=NPV$.
  \item Debt schedule: $V_{L,t}=$ remaining FCF at the WACC; $D_t=\frac DVV_{L,t}$;
        $ITS_t=\tau r_DD_{t-1}$; discount at $r_U$; check $V_L=V_U+PV(ITS)$ and
        $NPV_L=NPV_U+PV(ITS)$.
\end{enumerate}
\watch the shield at $t$ uses debt at $t-1$; debt goes to zero at the end of the project; owners'
value equals the NPV because outside investors get exactly their required return (zero NPV for them).

\blk{Last year's Q2 solved: the logic chain (60 marks; full solution in the mock's solutions PDF)}
\textbf{Data}: FCF $-750$ then $250$ for 5 years; comparable $D/E=0.1$, safe debt, $\beta_E=1.65$;
$r_f=5\%$, MRP $=10\%$, $\tau=30\%$.
\begin{center}\small
\begin{tabularx}{\linewidth}{@{}>{\raggedright\arraybackslash}p{1.1cm}>{\raggedright\arraybackslash}p{6.6cm}Y@{}}
  \toprule
  Part & Logic & Result \\ \midrule
  (a) & Constant $D/E$ $\Rightarrow$ shields move with firm value $\Rightarrow$ $\beta_{ITS}=\beta_U$,
        so $\beta_A=\beta_U$. Safe debt: $\beta_A=\frac EV\beta_E=\frac{1}{1.1}(1.65)$. CAPM for each
        claim; check $r_E$ with MM II. &
        $\beta_A=\beta_U=\beta_{ITS}=1.5$, $\beta_D=0$; $r_A=r_U=r_{ITS}=20\%$, $r_D=5\%$,
        $r_E=21.5\%$ \\[2pt]
  (b) & All equity: $r_A=r_E=r_U$. Annuity at 20\%: $250\times2.9906$. Negative NPV $\Rightarrow$
        outsiders would need $750/747.65>100\%$ of the firm. &
        $V_U=747.65$, $NPV=-2.35$; no financing \\[2pt]
  (c) & Constant $D/E$ $\Rightarrow$ MM II without $(1-\tau)$: $r_E=r_U+\frac DE(r_U-r_D)$;
        or $\beta_E=1.5(1.5)$. Larger: same asset risk on a smaller equity base. &
        $r_E=27.5\%$ ($\beta_E=2.25$) $>20\%$ \\[2pt]
  (d) & $D/V=1/3$: $r_{WACC}=\frac23(27.5\%)+\frac13(5\%)(0.7)$ $=r_U-\tau r_D\frac DV$.
        Annuity at 19.5\%: $250\times3.0238$. Reason: interest tax shields, not lower risk. &
        $r_{WACC}=19.5\%$, $V_L=755.95$, $NPV=+5.95$; yes \\[2pt]
  (e) & $D_0=\frac13V_L$, $E_0=\frac23V_L$. Outside equity puts in $750-D_0$ and buys
        $(750-D_0)/E_0$ of the equity at a fair price; owners keep the rest. &
        $D_0=251.98$, $E_0=503.97$; outsiders $98.82\%$; owners $1.18\%$ worth $5.95=NPV$ (equal) \\[2pt]
  (f) & $V_{L,t}=250\times$annuity$(19.5\%,5-t)$, $D_t=\frac13V_{L,t}$; $ITS_t=0.3(0.05)D_{t-1}$;
        discount at $r_U=20\%$. Link: $V_L=V_U+PV(ITS)$ and $NPV_{(d)}=NPV_{(b)}+PV(ITS)$. &
        $D_{1..5}=217.79, 176.92, 128.09, 69.74, 0$; $ITS_{1..5}=3.78, 3.27, 2.65, 1.92, 1.05$;
        $PV(ITS)=8.30=755.95-747.65$ \\ \bottomrule
\end{tabularx}
\end{center}
\watch the three ideas that carry the 60 marks: (i) constant $D/E$ $\Rightarrow$ $r_{ITS}=r_U$ (no
$(1-\tau)$ in MM II); (ii) negative all-equity NPV turns positive only through the tax shield;
(iii) outsiders earn zero NPV, so the owners' stake equals the NPV.

\recipe{WACC method vs.\ APV}
\textbf{WACC}: unlever comparables $\to$ $r_U$ $\to$ relever at target $D/E$ $\to$ $r_E$ $\to$
after-tax WACC $\to$ discount unlevered FCF. Valid only with a constant $D/E$.
\textbf{APV}: $V_U$ at $r_U$ $+$ $PV(ITS)$ at its own rate $-$ $PV(\text{distress})$ $+$ other side
effects. Best when the debt \emph{amount} is known (fixed loan, repayment schedule).
\watch never put the tax shield in the FCF \emph{and} in the WACC; the WACC weights are market values;
use the marginal tax rate.

\recipe{Risky debt and bankruptcy costs in states (PS2 Ex.\,3, Grinblatt--Titman 16.2)}
\begin{enumerate}
  \item Firm value: $V=\E[X]/(1+r_A)$ (or state prices).
  \item Creditors' payoff: $F$ if solvent; in default $X_s$ minus the bankruptcy cost (e.g.\ $(1-c)X_s$).
  \item Set $F$ so that creditors break even: $D=\E[\text{payoff}]/(1+r_D)$.
  \item Equity $=\E[\max(X-F,0)]$ discounted; compare with the no-cost case: the drop in equity value
        is the PV of bankruptcy costs.
\end{enumerate}
\watch the face value rises to compensate creditors; equity absorbs the whole cost ex ante.

\recipe{Static trade-off with a uniform distribution}
$X\sim U[\underline X,\bar X]$: $1-H(K)=\frac{\bar X-K}{\bar X-\underline X}$,
$h=\frac{1}{\bar X-\underline X}$ $\Rightarrow$ $K^*=\bar X-C/\tau$ if interior; else corner
$K^*=\underline X$ (safe debt only).
\watch check the corner: if $C\ge\tau(\bar X-\underline X)$ the firm takes only safe debt.

\recipe{Payout questions (Berk 17.6--17.13)}
Perfect market: dividend $\Rightarrow$ price falls by the dividend per share, wealth unchanged;
repurchase $\Rightarrow$ price unchanged, fewer shares. With taxes, compare the \emph{after-tax} value
investors get under each policy: corporate tax on retained interest $\Rightarrow$ paying out raises
value by $\tau_c\times$cash (perpetual); a dividend tax paid either way $\Rightarrow$ no change;
ex-dividend drop $=Div(1-\tau_d^*)$ with $\tau_d^*=\frac{\tau_d-\tau_g}{1-\tau_g}$.
\watch ``will this benefit investors?'' is a question about total wealth (cash $+$ shares), not the
share price.

% =============================================================================
\section*{4. Ready-made ``explain'' answers}
\begin{multicols}{2}\small
\textbf{Why is $r_E$ higher with more debt (even safe debt)?} The same asset risk is borne by a smaller
equity claim; debt's payments are fixed, so all variation falls on equity; $\beta_E$ rises and CAPM asks
a higher return. It is compensation for risk, not a cost increase.

\textbf{Why does the NPV rise with debt (taxes)?} Operating risk is unchanged; interest is deductible, so
the government takes less; the WACC falls below $r_U$ by $\tau r_DD/V$; equivalently NPV rises by
$PV(ITS)$.

\textbf{Why doesn't debt change value without taxes?} MM I: debt and equity only split the same cash
flow; investors can create or undo leverage on their own (homemade leverage), so they will not pay for
it.

\textbf{``Debt is cheaper than equity, so borrow more''?} Wrong: borrowing raises $r_E$ by exactly
enough to keep the WACC at $r_A$ (MM II); with taxes the gain is the shield, not ``cheap'' debt.

\textbf{Who bears bankruptcy costs?} Shareholders, ex ante: creditors price expected costs into the
amount they lend, so equity raises less.

\textbf{Why is the owners' stake equal to the NPV?} Outside investors pay a fair price (zero NPV);
the surplus, including the tax subsidy, stays with the owners.

\textbf{Why discount FCF and not earnings?} Earnings include non-cash items (depreciation) and ignore
investment and NWC; investors can only consume cash.

\textbf{Why unlever comparables?} Their equity betas mix business risk with their own leverage; only
the asset beta measures the business risk the project shares.

\textbf{Why does the price jump at the announcement, not at the repurchase?} Markets capitalise the PV
of the tax shields as soon as the permanent change is known; the repurchase itself is at a fair price.

\textbf{Why don't firms use 100\% debt?} Expected distress costs (direct and indirect), personal
taxes on interest, and agency costs offset the tax shield at the margin.

\textbf{Does a debt-financed dividend make the owner better off (MM)?} The owner gets cash today, but
equity falls by the same amount; firm value is unchanged, so total wealth is unchanged.

\textbf{Why can EPS rise but the price stay the same?} Higher EPS comes with higher risk; the required
return rises, so the P/E falls and the price does not move.
\end{multicols}

% =============================================================================
\section*{5. True/False strategy}
\begin{itemize}
  \item \textbf{Absolute words} (``always'', ``never'', ``any'', ``regardless'') are usually the trap:
        look for one counter-case (taxes, risky debt, distress, a different debt policy).
  \item \textbf{Hidden world}: decide first whether the statement lives in MM without taxes, with
        corporate taxes, or with distress costs. Many statements are true in one and false in another.
  \item \textbf{Converses}: ``positive NPV $\Rightarrow$ IRR $>r$'' may hold; the converse fails with
        non-standard flows.
  \item \textbf{Value vs.\ return}: ``raises the expected return on equity'' is not ``raises value''.
        ``Raises EPS'' is not ``raises the price''.
  \item \textbf{Book vs.\ market, per share vs.\ total, expected vs.\ realised, promised yield vs.\
        expected return}: a statement that swaps one for the other is false.
  \item \textbf{Who gets it}: tax shields to shareholders; distress costs paid by shareholders ex ante;
        NPV to existing owners.
  \item \textbf{One claim per item}: if any part of a compound statement is wrong, it is false.
  \item \textbf{Penalty}: under a ``wrong answers cost half a right one'' rule (the MPE rule used in this
        project's lists), guessing pays only if you are more than one-third sure; otherwise leave blank.
\end{itemize}

% =============================================================================
\section*{6. True/False bank (75 items)}
\begin{multicols}{2}\small
\blk{Topic 1}
\begin{enumerate}[label=\arabic*.,leftmargin=1.6em,itemsep=2pt]
  \item A firm with negative book equity is insolvent. \F Buybacks or accumulated losses can push
        book equity below zero while market equity is large (Starbucks 2020).
  \item Depreciation is a source of cash. \F It is non-cash; only the tax it saves, $\tau\,$Dep,
        affects cash.
  \item An increase in accounts receivable makes operating cash flow lower than net income, other
        things equal. \T Sales booked but not collected.
  \item An increase in accounts payable lowers cash flow from operations. \F It raises CFO: suppliers
        finance the firm.
  \item If capex equals depreciation, net income equals free cash flow. \F It also needs $\Delta
        NWC=0$ and no interest (NI is after interest; FCF is unlevered).
  \item Enterprise value equals market capitalisation plus total liabilities. \F Market cap plus
        (interest-bearing) debt minus cash.
  \item Higher leverage always raises ROE. \F Only if $ROIC>r_D(1-\tau)$; it always raises ROE's
        volatility.
  \item Two firms with the same ROE are equally profitable in their operations. \F DuPont: margin,
        turnover and leverage can differ; ROE ignores risk and timing.
  \item ROIC does not depend on how the firm is financed. \T $EBIT(1-\tau)$ over total invested
        capital, whatever its split.
  \item Book and market debt-to-equity ratios give the same leverage. \F Book equity is historical;
        market equity is forward-looking.
  \item Under accrual accounting, revenue is recognised when cash is received. \F When it is earned.
  \item Higher fixed costs make operating income more sensitive to sales. \T $DOL=1+F/OI$.
  \item A growing firm can report positive net income and negative operating cash flow. \T Growth
        ties up cash in receivables and inventory.
  \item Interest expense is deducted when computing unlevered FCF. \F Financing stays out of FCF.
  \item An increase in deferred tax liabilities is added back from net income to cash flow. \T Taxes
        expensed but not yet paid (PS1 Ex.\,1b).
  \item The quick ratio cannot exceed the current ratio. \T It removes inventory from the numerator.
  \item A high market-to-book ratio means the stock is overvalued. \F It usually reflects growth and
        intangibles missing from the balance sheet.
  \item Excess cash belongs in NWC when computing FCF. \F NWC counts operating items only.
\end{enumerate}

\blk{Topic 2A}
\begin{enumerate}[label=\arabic*.,start=19,leftmargin=1.6em,itemsep=2pt]
  \item A project should be discounted at the firm's WACC. \F Only if it has the firm's risk and
        financing; use the project's own risk.
  \item A levered comparable's equity beta is the right beta for an all-equity project in the same
        business. \F Unlever it first.
  \item If debt is risk-free, issuing it does not change the risk of equity. \F The same asset risk
        falls on less equity: $\beta_E$ rises.
  \item Sunk costs belong in the project's cash flows. \F They do not change with the decision.
  \item The value of an asset the firm already owns and would use in the project is a cost of the
        project. \T Opportunity cost.
  \item Allocated overhead that is incurred anyway is an incremental cost. \F Not incremental.
  \item A project with a higher IRR always has a higher NPV. \F Scale and timing differ.
  \item For any cash-flow pattern, accept iff $IRR>r$. \F Borrowing-type flows and multiple IRRs
        break the rule.
  \item In a frictionless market, a project's NPV accrues to existing shareholders whether it is
        financed with equity or safe debt. \T The financing is zero-NPV.
  \item State prices are the probabilities of the states. \F They also reflect discounting and risk
        aversion; they sum to $1/(1+r_f)$.
  \item The state prices sum to $1/(1+r_f)$. \T A portfolio of all claims is a risk-free bond.
  \item No arbitrage requires all state prices to be strictly positive. \T Otherwise a claim with
        a positive payoff would be free.
  \item In a perfect market, borrowing to pay a dividend raises firm value. \F Value is unchanged;
        equity falls by the dividend (PS2 1f).
  \item Without taxes, more debt lowers the WACC because debt is cheaper than equity. \F $r_E$ rises
        to offset it; WACC $=r_A$.
  \item $r_E$ rises linearly with $D/E$ even when debt becomes risky. \F Linear only while debt is
        safe; flatter once $r_D$ rises.
  \item Equity is a call on the firm's assets with strike equal to the face value of debt. \T
        $\max(A-F,0)$.
  \item Risky debt equals a risk-free bond plus a put on the assets. \F Minus a put: creditors have
        written it.
  \item The yield to maturity of risky debt is its expected return. \F $y>r_D$ when default is
        possible.
  \item MM's arbitrage proof needs investors to borrow on the same terms as firms. \T Homemade
        leverage.
  \item If a recapitalisation raises the expected return on equity, shareholders are better off. \F
        It compensates for added risk.
  \item In a perfect market, leverage that raises EPS raises the share price. \F The price is
        unchanged (PS2 Ex.\,4).
  \item If a debt-financed repurchase raises EPS in a perfect market, the P/E falls. \T Same price,
        higher (riskier) EPS.
  \item The beta of a portfolio is the value-weighted average of the betas of its parts. \T This is
        why $\beta_A=\frac EV\beta_E+\frac DV\beta_D$.
  \item $r_A=\frac EVr_E+\frac DVr_D$ holds only for expected returns, not realised ones. \F It holds
        for realised returns too (the slides use this to get betas).
  \item Fisher separation: investment depends on the owners' time preferences. \F The market rate
        decides; owners borrow or lend privately.
  \item A positive-NPV project can survive in a competitive market because it cannot be scaled. \T
        ``Arbitrage across the firm's boundary''.
  \item Real cash flows may be discounted at nominal rates. \F Real with real, nominal with nominal.
  \item Working capital recovered at the end of the project has no effect on NPV. \F It is tied up
        for years: it costs the time value of money.
  \item The after-tax salvage value is $SV(1-\tau)$. \F It is $SV-\tau(SV-BV)$.
\end{enumerate}

\blk{Topic 2B}
\begin{enumerate}[label=\arabic*.,start=48,leftmargin=1.6em,itemsep=2pt]
  \item With corporate taxes, $V_L-V_U=\tau D$ for any debt policy. \F Only for permanent fixed debt.
  \item The interest tax shield benefits creditors. \F Debt is fairly priced; the gain goes to
        shareholders.
  \item After the announcement of a permanent debt-financed buyback, the price rises at once, not on
        the repurchase date. \T The shield is capitalised on the news.
  \item With taxes and debt, the after-tax WACC is below $r_U$. \T By $\tau r_D D/V$ (constant $D/E$).
  \item Unlevered FCF should include the interest tax shield. \F It enters via WACC or APV.
  \item WACC and APV give the same value when $D/E$ is constant. \T Both discount the shields at $r_U$.
  \item APV is the easier method when the debt amount follows a known schedule. \T WACC needs a
        constant $D/E$.
  \item With a constant $D/E$, the tax shields are as risky as the debt. \F As risky as the
        unlevered assets: $r_{ITS}=r_U$.
  \item With permanent fixed debt, the shields are discounted at $r_D$. \T Their amount is fixed like
        the debt's payments.
  \item With constant $D/E$, $r_E=r_U+\frac DE(1-\tau)(r_U-r_D)$. \F That is the permanent-debt
        formula; with constant $D/E$ there is no $(1-\tau)$.
  \item MM with corporate taxes alone predicts nearly all-debt firms. \T $V_L/V_U=1+\tau D/V_U$.
  \item Personal taxes on interest income can shrink or cancel the tax advantage of debt. \T Miller.
  \item Default by itself violates MM. \F Default is a transfer of control; only its costs matter.
  \item Bankruptcy costs are borne by creditors, who receive less in default. \F Creditors price them
        in; shareholders bear them ex ante (PS2 Ex.\,3).
  \item Indirect distress costs arise only after a bankruptcy filing. \F Customers, suppliers and
        workers react before (Hortaçsu et al.).
  \item Firms whose products need long-term service should borrow less under the trade-off theory. \T
        Their indirect distress costs are high.
  \item In the static trade-off, optimal debt rises with $\tau$ and falls with bankruptcy costs. \T
        From $\tau[1-H(K^*)]=C\,h(K^*)$.
  \item At the optimal debt level, the PV of the tax shield is maximised. \F Firm value is: marginal
        shield $=$ marginal expected distress cost.
  \item More equity raises a bank's WACC because equity is costlier than debt. \F MM II: $r_E$ falls
        with less leverage; only the lost tax shield is a (private) cost.
  \item More equity lowers a bank's expected ROE, so its shareholders lose. \F ROE's risk falls too;
        the required return falls.
  \item In a perfect market, switching from dividends to repurchases raises value. \F Payout is
        irrelevant.
  \item Dividend irrelevance requires the investment policy to be held fixed. \T Otherwise payout
        changes investment.
  \item In a perfect market, the stock price falls by the dividend on the ex-dividend date. \T Value
        leaves the firm.
  \item In a perfect market, an open-market repurchase raises the share price. \F Price unchanged,
        fewer shares.
  \item If dividends are taxed more heavily than capital gains, repurchases are tax-preferred. \T Lower
        effective tax.
  \item Holding excess cash is costly when the firm pays corporate tax on its interest income. \T
        Investors could earn it untaxed (Berk 17.12).
  \item With risky debt, the tax shield is computed with the expected return $r_D$. \F With the
        coupon actually paid, $\bar r_D$.
  \item With taxes and riskless debt, leverage still raises the equity beta. \T The equity bears the
        asset risk on a smaller base.
\end{enumerate}
\end{multicols}

\end{document}
